\section{Discussion}
\label{sec:discussion}

\subsection{Why the Hypothesis Fails}
\label{sec:why_fails}

The artificial outsider hypothesis rests on three premises: (1) LLMs share evaluation biases that make them prefer homogeneous text; (2) text that deviates from the LLM-preferred distribution is more novel; and (3) novelty translates to higher human interest.
Our results show that premise (3) fails.
While the LLM judges do agree closely with each other ($\rho \approx 0.75$), lending some support to premise (1), deviation from LLM preferences does not predict higher human interest.
LLM-disliked text is simply lower quality, and humans also find lower-quality text less interesting.

\subsection{The Quality Confound}
\label{sec:quality_confound}

The most important finding of this work is the quality confound.
The raw correlation of $\rho = +0.48$ between LLM consensus and human interestingness drops to $\rho \approx 0.00$ after controlling for overall text quality.
This has three implications.
First, LLMs and humans agree almost perfectly on what constitutes good versus bad text.
Second, beyond quality, LLMs have no additional bias for or against interesting content.
Third, the outsider hypothesis conflates ``quality outlier'' with ``interestingness outlier''---these are different things, and the former dominates.

This finding also constrains how to interpret prior work on LLM evaluation biases.
While LLMs exhibit self-preference~\citep{panickssery2024llm} and perplexity preference~\citep{wataoka2024self}, these biases appear to be largely \emph{aligned} with human quality judgments rather than opposed to them.
The biases are real, but they do not cause LLMs to systematically miss interesting content.

\subsection{When Might the Hypothesis Hold?}
\label{sec:when_hold}

Our results do not rule out the outsider hypothesis in all settings.
Three conditions not tested here could produce different results.

\para{Modern high-quality LLMs.}
The \hanna dataset uses generators from the GPT-2 era (2019--2022), producing text with high quality variance---from near-gibberish to coherent prose.
For modern models where all outputs are competent, the quality confound is weaker, and an interestingness signal might emerge.

\para{Specific creative domains.}
Creative writing at the frontier---poetry, literary fiction, experimental prose---may reward rule-breaking in ways that standard quality metrics do not capture.
In these domains, ``interesting'' may genuinely mean ``deviating from convention'' rather than ``well-executed.''

\para{Different operationalization.}
We proxy interestingness through Surprise and Engagement.
A more targeted measure---perhaps focused on novelty of ideas or unexpectedness of narrative structure---might capture aspects of interestingness that LLMs do systematically undervalue.

\subsection{Limitations}
\label{sec:limitations}

\para{Dataset age.}
\hanna's generators are from the GPT-2 era.
Text quality varies from near-gibberish to coherent, and the quality confound may dominate precisely because quality variance is so large.
Our findings may not generalize to evaluations of modern, uniformly competent LLM outputs.

\para{LLM judge diversity.}
All three judges are OpenAI models.
Using models from different families (Claude, Gemini, Llama) could reveal more meaningful disagreement patterns.
The high inter-judge agreement we observe ($\rho = 0.71$--$0.77$) may partly reflect shared training data or alignment procedures.

\para{Interestingness proxy.}
We use Surprise and Engagement from \hanna, not a direct ``interestingness'' rating.
These dimensions may not fully capture the concept, particularly aspects like intellectual stimulation or memorability.

\para{Domain specificity.}
Our results come from short creative fiction (\hanna) and general chatbot conversations (\arena).
Other domains---research ideas, essays, journalism---may produce different human-LLM alignment patterns.

\para{Arena sample size.}
We evaluate only 200 Arena battles, limiting statistical power for subgroup analyses.
